% Chapter 1. Ported from docs/chapter1-draft.md (15 Oct 2026), updated 13 Nov 2026 for the island landscape.
\section{INTRODUCTION}

\subsection{Background of the Study}

Negros Island is the centre of Philippine sugarcane, with about 57\% of the national sugarcane area \citep{PCAARRDnd}. The crop supports a large and mostly smallholder base, with about 88,000 sugarcane farmers nationally, of whom about 84\% hold less than 5 ha \citep{PCAARRDnd}. An earlier count gives 78,276 farmers and 81.5\% below 5 ha for 2013--2014 \citep{Briones2020}. Shocks to the crop therefore turn quickly into rural income loss.

Since 2022 the industry has faced a new threat, the red-striped soft scale insect (RSSI), \emph{Pulvinaria tenuivalvata} (Newstead) (Hemiptera: Coccidae). RSSI is a sap-sucking soft scale that excretes honeydew, which promotes sooty mould and interferes with photosynthesis \citep{Guerrero2026}. Earlier studies reviewed by \citet{Guerrero2026} report that this feeding reduces cane yield and sucrose content, although no such study has been done in the Philippines \citep{Guerrero2024}. The insect reproduces by thelytokous parthenogenesis, so populations are all female and can build up without mating \citep{Guerrero2026}. Under Philippine laboratory conditions a female lays about 197 eggs, and the cycle from egg to death lasts about 63 days \citep{Guerrero2026}. RSSI was first observed in the Philippines in Bacolor, Pampanga, in June 2022 \citep{Navasero2023}. In Negros it was first noted in northern Negros Occidental in late March 2025 and was declared an infestation on 22 May 2025, when it had spread over 87 ha \citep{ManilaBulletin2025}. By 2 October 2026 the Sugar Regulatory Administration (SRA) had validated 72,103 ha of RSSI-affected sugarcane on Negros Island, about 30\% of the 239,216 ha planted in crop year 2025--2026 \citep{SRA2026FOI}.

Negros Occidental was placed under a state of calamity in July 2025 \citep{Masculino2025} and again on 30 June 2026 \citep{Bayoran2026}, and Negros Oriental followed on 20 July 2026 \citep{Partlow2026}. The response so far has relied on local government units and farmers to deliver control \citep{Partlow2026,Gomez2026}. It includes chemical and drone spraying, biological control with entomopathogenic fungi, and emergency cash assistance \citep{ChiniMandi2026}, quick response funds released under the calamity declarations \citep{Partlow2026}, and the propagation of a variety reported to be more resistant \citep{Gomez2026}. The provincial government of Negros Occidental has declined to fund chemical spraying and has opted for biological control, although it has not banned chemicals \citep{Bayoran2026}. Its Provincial Biosecurity Task Force applies the entomopathogenic fungi \emph{Metarhizium} and \emph{Beauveria}, by truck-mounted mist sprayer where a farm track reaches the field and through farmers' backpack sprayers elsewhere \citep{PGNO2026}, and the province is propagating that variety \citep{Gomez2026}. The task force also surveys infestation in the field. As of 8 November 2026 it had surveyed 16,505 farmlots in all 31 local government units it covers in Negros Occidental: 50\% were rated low, 29\% moderate, and 15\% high, and no RSSI was found in 6\% \citep{PGNO2026}. The survey is a dated snapshot whose coverage has grown over time, so it does not show how infestation has changed. The choice between biological, chemical, and mixed control is being made in real time and at scale, without a published or operationally deployed decision-support tool to compare control strategies before deployment.

\subsection{Statement of the Problem}

The RSSI outbreak in Negros is severe, spreading, and managed reactively, and the central management question, which control strategy to pursue, is being decided without an established quantitative basis. The system is data-poor. The pest's effect on cane yield and sucrose content has not been studied in the Philippines \citep{Guerrero2024}, and the absence of region-specific biological data limits the development of monitoring systems and economic threshold levels \citep{Guerrero2026}. Infestation levels are now defined by insect counts per leaf \citep{PGNO2026}, but no study links those counts to yield loss, so there is still no damage-based threshold for deciding when to act.

Agent-based modelling can represent the pest, the crop, and the farmer as interacting agents and can be used to explore management strategies in silico before they are deployed in the field. A systematic review of the agent-based pest-management literature (Review of Literature) found that, of 30 reviewed models, none models a scale insect or a sugarcane-RSSI system, and none is validated on Philippine field data. No agent-based model for comparing RSSI control strategies on sugarcane has been reported. The problem this study addresses is to build such a model when many of the quantities it needs, such as spread between fields and yield loss, have not been measured, and to compare control strategies while stating clearly which inputs are known and which are assumed.

\subsection{Objectives of the Study}

The general objective is to develop an agent-based model of red-striped soft scale insect management on sugarcane in Negros Island as a tool for exploring pest-management strategies. The specific objectives are as follows.

\begin{enumerate}
	\item To characterise the RSSI-sugarcane-farmer system in Negros from literature and field reports, including pest biology, crop response, control options, and the management context.
	\item To formalise a conceptual model of pest, crop, farmer, and natural-enemy interactions suitable for agent-based implementation.
	\item To identify data sources and to state explicitly what can and cannot be parameterised, treating unsupported quantities as assumptions to be tested.
	\item To implement an initial agent-based model in GAMA.
	\item To evaluate model behaviour through exploratory simulations comparing control strategies (chemical, biological, hybrid, and none) on cumulative pest load, the extent of pest spread over the simulated area, and the number of control applications, with cane yield loss derived from an assumed damage function, allowing the efficacy of chemical control to decline with repeated use so that the strategies are compared over sustained use rather than a single application.
\end{enumerate}

\subsection{Significance of the Study}

The management decision facing the province, biological versus chemical versus hybrid control, is currently made without a way to test strategies before committing resources. A model that represents the system and lets strategies be compared in silico offers a low-cost testbed for that decision. Because chemical efficacy is allowed to fall with repeated use, the comparison also speaks to how each strategy holds up under sustained use, which is the crux of the biological-versus-chemical debate, rather than only to a single season.

The study is also useful for its method. By recording which parameters are supported by data and which are not, it documents what a data-poor, emerging-pest system can and cannot yet support in a model. This record is useful to the SRA and to later researchers.

\subsection{Scope and Limitations}

This study develops a conceptual model and an initial, exploratory agent-based model of RSSI management on sugarcane in Negros. The model is intended for comparing control strategies. Its outputs are relative comparisons between strategies, not calibrated field predictions.

Several quantities cannot be parameterised from available data. Pest life-history parameters come from a single laboratory temperature regime \citep{Guerrero2026}, so the model uses no development function that depends on temperature. No dose-response relates scale density to yield or sugar loss, so the damage function takes an assumed form. Crawler dispersal, stage-specific field mortality, and per-hectare control costs could not be obtained from verifiable sources and are treated as assumptions swept across plausible ranges. Pesticide resistance is not modelled as an evolutionary process. Instead, the efficacy of chemical control declines with cumulative use at an assumed rate, so that the long-run disadvantage of repeated spraying is represented without genotype-level data, which were not found for RSSI. Because these quantities are uncalibrated, the model is assessed with a hindcast against one verified Negros figure, not by statistical fitting.

The study is limited to RSSI on sugarcane in Negros and to the comparison of chemical, biological, hybrid, and no-control strategies. The landscape covers all of Negros Island at 1-ha resolution and is built from a satellite sugarcane map for 2019--2022 \citep{DiTommaso2024}. That map predates the outbreak and overstates the area planted in 2025--2026, the SRA reply included no infestation map, because mapping is still ongoing, and no key linking municipal records to mill districts \citep{SRA2026FOI}, and the provincial survey records are viewable but not released as data \citep{PGNO2026}. The model therefore starts from 87 infested cells placed at random within the municipalities first reported as affected and is not used to say where the pest will spread next. Re-invasion from outside the island is outside the scope of the study. Reported losses of up to 50\% of cane yield and sugar content \citep{ChiniMandi2026} are treated as claims from news sources, not as established measurements, because no primary study is cited for them.

\subsection{Dataset}

The study does not use a single measured dataset. It assembles parameters and context from published literature, government and industry sources, open geospatial data, and news reports. Each item is recorded with its source and a status of real, analogue, or assumption in the study's parameter table.

Pest biology, the best-supported component, comes from \citet{Guerrero2026}, which gives stage durations, fecundity, reproduction mode, and longevity for a Philippine laboratory population. Control efficacy has Philippine anchors. \emph{Metarhizium anisopliae} caused 86.67\% to 91.11\% mortality at 15 days in vitro \citep{Buletin2026}. Natural parasitism of teneral and young adults was 34.85\% in a screenhouse and 66.8\% in the field in Luzon \citep{Guerrero2024}. The cropping calendar comes from \citet{MojicaSevilla2025}: the crop year runs from October to September, and planting runs from October to May. Negros cane yield and planted area by mill district for crop years 2020--2021 to 2025--2026, sugar recovery for 2025--2026, and the validated RSSI-affected area by municipality as of 2 October 2026, come from the SRA reply to a freedom-of-\allowbreak information request \citep{SRA2026FOI}. Across the 12 Negros mill districts, cane yield ranged from 49.32 to 71.64 tonnes per hectare over crop years 2020--2021 to 2024--2025, and Negros as a whole recorded 59.48 tonnes per hectare and a sugar recovery of 1.61 LKg per tonne of cane in 2025--2026 \citep{SRA2026FOI}. The definitions of infestation Levels 1 to 3 and the current severity of infestation in Negros Occidental come from the provincial monitoring dashboard \citep{PGNO2026}, an official operational source that is dated and subject to revision. The landscape uses the 10-m sugarcane map of \citet{DiTommaso2024}, municipal boundaries from the OCHA common operational dataset \citep{OCHA2025}, and roads from OpenStreetMap \citep{OSM2026}. Sugar prices and the farm-size distribution give context but are not model inputs, because farm income is not modelled.

The dataset has notable gaps. No economic injury level, economic threshold, or damage function was found for RSSI, and the lack of region-specific biological data limits their development \citep{Guerrero2026}. The provincial dashboard describes damage that worsens with severity, with yield able to fall at high severity, but it gives no measured relation \citep{PGNO2026}. Per-hectare control costs, the farmgate cane price, and labour costs were not found in verifiable sources. The SRA reports that its mapping is still ongoing \citep{SRA2026FOI}, and the provincial survey is a single dated snapshot, so no multi-date infestation series was found. Climate is held out of the model, so development runs at the laboratory rate of \citet{Guerrero2026}. No resistance or efficacy-loss data were found for RSSI, so the decline in chemical efficacy with repeated use is set as an assumption. These gaps define the parameters that are handled as assumptions and tested by sensitivity analysis rather than presented as data.
